\section{Stages}
\label{sec:stages}

A \emph{legal stage of arity \(d\)}, or simply a \emph{stage}, written
\((r,b)\to(c,r)\), is given by integers with
\begin{equation}\label{eq:stage-conditions}
 r\ge1,\qquad d\ge2,\qquad c\ge1,\qquad (d+2)c\le db-r+1 .
\end{equation}
The last condition implies \(c<b\), as \(dc<(d+2)c\le db-r+1\le db\).
The stage
has \emph{parent} \(C(r,b)\), \emph{child length} \(c\) and next pair
\((c,r)\). We put \(h=b-c\ge1\) and \(b^+=db-(d-1)c=b+(d-1)h\), the
\emph{shift} and the \emph{quotient length}.
Lemma~\ref{lem:shifted-chain} splits \(C(r,b)B_X(d)\)
block-triangularly into \(d-1\) kernel parts \(C(r,c)\) and a quotient
part \(C(r,b^+)\). By Lemma~\ref{lem:reset}(c), the last condition
of~\eqref{eq:stage-conditions} says that \(b^+\ge3c+r-1\), so
Lemma~\ref{lem:reset} can crop the quotient to \(C(r,3c+r-1)\) and
reset it into two copies of \(C(r,c)\) and one of \(C(r,c)^{(23)}\).
Lemma~\ref{lem:stage-types} deletes all cross terms in a tensor power by
fixing class counts, and Proposition~\ref{prop:stage},
Corollary~\ref{cor:exact-stage} and Proposition~\ref{prop:chain-counts}
turn this into exact relations and compose them.

\subsection{The shifted sum}
\label{subsec:shifted-chain}

Let \(\mathbb F\) be a field, and let \(\mathbb F[\xi]_{<\ell}\) be the
polynomials in \(\xi\) of degree less than \(\ell\), with basis
\(\xi^0,\ldots,\xi^{\ell-1}\). Fix \(d\ge2\) and \(h\ge1\). For
\(\ell\ge h+1\) let \(W_\ell=(\mathbb F[\xi]_{<\ell})^d\), with
components \(j=0,\ldots,d-1\) and monomial basis \(\xi^me_j\), and define
the \emph{shifted sum}
\[
 W_\ell\to\mathbb F[\xi]_{<\ell+(d-1)h},\qquad
 (g_0,\ldots,g_{d-1})\longmapsto\sum_{j=0}^{d-1}\xi^{jh}g_j .
\]
For \(0\le j\le d-2\) and \(g\in\mathbb F[\xi]_{<\ell-h}\), the
\emph{kernel vector} \(\kappa_j(g)\in W_\ell\) has component \(j\) equal
to \(-\xi^hg\), component \(j+1\) equal to \(g\), and all others zero.
For \(0\le\gamma<\ell+(d-1)h\), let \(j_\ell(\gamma)\) be the least \(j\)
with \(\gamma\le jh+\ell-1\), and let
\(\lambda_\ell(\gamma)=\xi^{\gamma-j_\ell(\gamma)h}e_{j_\ell(\gamma)}\),
the \emph{lift} of \(\xi^\gamma\). It is a monomial of \(W_\ell\), the
one in the first component that contains the exponent \(\gamma\): if
\(j=j_\ell(\gamma)\ge1\), then \(\gamma>(j-1)h+\ell-1\ge jh\).

For integers \(r\ge1\), \(d\ge2\) and \(1\le c<b\), with \(h=b-c\), the
tensor
\begin{equation}\label{eq:stage-U}
 U=C(r,b)\,B_X(d)=\sum_{\alpha=0}^{r-1}\sum_{\beta=0}^{b-1}\sum_{j=0}^{d-1}
 x_\alpha\,y_{j,\beta}\,z_{j,\alpha+\beta}
\end{equation}
multiplies \(f\in\mathbb F[\xi]_{<r}\) by \(g\in W_b\) componentwise: the
variable \(z_{j,\gamma}\) reads the coefficient of \(\xi^\gamma e_j\) in
\(fg\in W_{b+r-1}\).

\begin{lemma}[Shifted sum]\label{lem:shifted-chain}
Let \(d\ge2\), \(h\ge1\) and \(\ell\ge h+1\).
\begin{enumerate}
\item[(a)] The lifts \(\lambda_\ell(\gamma)\), \(0\le\gamma<\ell+(d-1)h\),
and the kernel vectors \(\kappa_j(\xi^i)\), \(0\le j\le d-2\),
\(0\le i<\ell-h\), form a basis of \(W_\ell\), the \emph{shifted basis}.
Its matrix has entries in \(\{0,\pm1\}\) and determinant \(\pm1\), and its
inverse has entries in \(\{0,1\}\): for \(0\le m<\ell\) and \(0\le j<d\),
\begin{equation}\label{eq:stage-inverse}
 \xi^me_j=\lambda_\ell(jh+m)
 +\sum_{j'=j_\ell(jh+m)}^{j-1}\kappa_{j'}\bigl(\xi^{m+(j-1-j')h}\bigr).
\end{equation}
\item[(b)] The shifted sum is surjective, and its kernel is the direct
sum of the \(d-1\) subspaces \(\kappa_j(\mathbb F[\xi]_{<\ell-h})\).
\item[(c)] Let \(r\ge1\) and \(1\le c<b\), with \(h=b-c\). Use the shifted
basis of \(W_b\) on leg~\(2\) of \(U\), the coordinates dual to the
shifted basis of \(W_{b+r-1}\) on leg~\(3\), and the monomial basis on
leg~\(1\). Call the new coordinates \(y^{\mathrm q}_\beta\) and
\(y^{(j)}_i\) (for \(\lambda_b(\beta)\) and \(\kappa_j(\xi^i)\)), and
\(z^{\mathrm q}_\gamma\) and \(z^{(j)}_i\) (for
\(\lambda_{b+r-1}(\gamma)\) and \(\kappa_j(\xi^i)\)). Then \(U\) becomes
\begin{equation}\label{eq:stage-tensor}
 U'=\sum_{\alpha=0}^{r-1}x_\alpha\Biggl(
 \sum_{\beta=0}^{b^+-1}y^{\mathrm q}_\beta z^{\mathrm q}_{\beta+\alpha}
 +\sum_{j=0}^{d-2}\sum_{i=0}^{c-1}y^{(j)}_iz^{(j)}_{i+\alpha}
 +\sum_{\beta=0}^{b^+-1}\ \sum_{j'=j_{b+r-1}(\beta+\alpha)}^{j_b(\beta)-1}
 y^{\mathrm q}_\beta z^{(j')}_{\beta+\alpha-(j'+1)h}\Biggr).
\end{equation}
The first sum is \(C(r,b^+)\), the second is \(d-1\) copies of \(C(r,c)\),
all sharing leg~\(1\), and the third, which joins quotient to kernel
coordinates, is nonzero if and only if \(r\ge2\). The substitutions are
integral with integral inverses, so \(U\restr U'\restr U\) over every
field.
\end{enumerate}
\end{lemma}

\begin{proof}
(a) For each \(\gamma\), the monomials \(m_j=\xi^{\gamma-jh}e_j\) of
\(W_\ell\) have \(j_0\le j\le j_1\), where \(j_0=j_\ell(\gamma)\) and
\(j_1=\min(d-1,\lfloor\gamma/h\rfloor)\); let \(W^{(\gamma)}\) be their
span. Every monomial \(\xi^me_j\) lies in \(W^{(m+jh)}\), so
\(W_\ell=\bigoplus_\gamma W^{(\gamma)}\). The lift \(\lambda_\ell(\gamma)\)
is \(m_{j_0}\), and \(\kappa_j(\xi^i)=\xi^ie_{j+1}-\xi^{i+h}e_j\) lies in
\(W^{(\gamma)}\) for \(\gamma=i+(j+1)h\), where the conditions
\(0\le i<\ell-h\) and \(j\le d-2\) read \(j_0\le j<j_1\). So the vectors
of~(a) in \(W^{(\gamma)}\) are \(m_{j_0}\) and \(m_{j+1}-m_j\),
\(j_0\le j<j_1\). Their matrix in the basis \(m_{j_0},\ldots,m_{j_1}\) is
unitriangular with entries \(0,\pm1\), and its inverse
\(m_j=m_{j_0}+\sum_{j'=j_0}^{j-1}(m_{j'+1}-m_{j'})\)
is~\eqref{eq:stage-inverse}. Hence (a) holds on every \(W^{(\gamma)}\),
and so on \(W_\ell\).

(b) The shifted sum maps \(\xi^me_j\) to \(\xi^{m+jh}\), so it maps the
lifts bijectively onto the monomials \(\xi^\gamma\) and is onto; and it
maps \(\kappa_j(g)\) to \(-\xi^{(j+1)h}g+\xi^{(j+1)h}g=0\). As the lifts
and kernel vectors form a basis by~(a), the kernel vectors form a basis of
the kernel.

(c) As in \citep[Section~4.1]{FinCon26}, substitute \(y=\Pi_by'\) and
\(z=\Pi_{b+r-1}^{-\top}z'\), where \(\Pi_\ell\) has the shifted basis as
columns; the coefficient of \(x_\alpha y'_{\mathbf u}z'_{\mathbf v}\) is
then the coordinate along \(\mathbf v\) of \(\xi^\alpha\mathbf u\). First,
\(\xi^\alpha\kappa_j(\xi^i)=\kappa_j(\xi^{i+\alpha})\) in \(W_{b+r-1}\),
since \(i+\alpha<c+r-1=(b+r-1)-h\). Second,
\(\xi^\alpha\lambda_b(\beta)=\xi^{\beta+\alpha-jh}e_j\) with
\(j=j_b(\beta)\), and by~\eqref{eq:stage-inverse} for \(W_{b+r-1}\) this
is \(\lambda_{b+r-1}(\beta+\alpha)\) plus
\(\kappa_{j'}(\xi^{\beta+\alpha-(j'+1)h})\) for
\(j_{b+r-1}(\beta+\alpha)\le j'<j_b(\beta)\). This
gives~\eqref{eq:stage-tensor}; in the first sum, \(\beta\) runs over
\([0,b^+)\) and \(\beta+\alpha\) over \([0,b^++r-1)\). If \(r=1\), then
\(\alpha=0\) and \(j_{b+r-1}=j_b\), so the third sum is empty. If
\(r\ge2\), then \(\beta=b<b^+\) and \(\alpha=0\) give \(j_b(b)=1\) and
\(j_{b+r-1}(b)=0\), so the third sum contains
\(x_0y^{\mathrm q}_bz^{(0)}_c\).
\end{proof}

\begin{example}\label{ex:shifted-chain}
Take \(d=2\), \(r=2\), \(b=3\) and \(c=1\), so that \(h=2\), \(b^+=5\)
and the shifted sum is \((g_0,g_1)\mapsto g_0+\xi^2g_1\). The shifted
bases are
\[
\begin{array}{c|l|l}
 &\text{lifts }\lambda_\ell(0),\lambda_\ell(1),\ldots&\text{kernel vectors }\kappa_0(\xi^i)\\ \hline
 W_3&e_0,\ \xi e_0,\ \xi^2e_0,\ \xi e_1,\ \xi^2e_1&e_1-\xi^2e_0\\
 W_4&e_0,\ \xi e_0,\ \xi^2e_0,\ \xi^3e_0,\ \xi^2e_1,\ \xi^3e_1&e_1-\xi^2e_0,\ \xi e_1-\xi^3e_0
\end{array}
\]
Multiplying a vector of the first row by \(1\) or \(\xi\) gives a vector
of the second row in every case but one: \(1\cdot\lambda_3(3)=\xi e_1\)
is \(\lambda_4(3)+\kappa_0(\xi)\) in \(W_4\). Hence \(U'\) consists of
\(C(2,5)\) in the coordinates \(x,y^{\mathrm q},z^{\mathrm q}\), of
\(C(2,1)\) in \(x,y^{(0)},z^{(0)}\), and of the single term
\(x_0y^{\mathrm q}_3z^{(0)}_1\).
\end{example}

\paragraph{Binary forms.}
Identify \(\xi^me_j\in W_\ell\) with
\(u^{\ell-1-m}v^m\,u'^{\,d-1-j}v'^{\,j}\in V_\ell\otimes V'_d\), where
\(V_\ell\) and \(V'_d\) are the binary forms with \(\ell\) and \(d\)
coefficients in \(u,v\) and in \(u',v'\) of \citep[Section~4.1]{FinCon26},
and identify \(\xi^\gamma\) with \(u^{\ell+(d-1)h-1-\gamma}v^\gamma\). Then
the shifted sum is the substitution \((u',v')=(u^h,v^h)\), and
\(\kappa_j(\xi^i)=\Delta_h\,u^{\ell-h-1-i}v^iu'^{\,d-2-j}v'^{\,j}\) with
\(\Delta_h=u^hv'-v^hu'\). For \(h=1\), the kernel
\(\Delta_1(V_{\ell-1}\otimes V'_{d-1})\) is the first step
\(\mathcal F^1(V_\ell\otimes V'_d)\) of the determinant filtration
of \citep[Section~4.1]{FinCon26}, defined there for \(\ell\ge d\). For
\(b\ge d\), the split of width \(d-1\) there refines \(\mathcal F^1\)
further, with quotients carrying \(C(r,b+d-1-2j)\), \(1\le j\le d-1\);
Lemma~\ref{lem:shifted-chain} keeps it whole, as \(d-1\) equal parts
\(C(r,c)\).

\subsection{Crop and reset}
\label{subsec:reset}

\begin{lemma}[Crop and reset]\label{lem:reset}
Let \(r,c\ge1\).
\begin{enumerate}
\item[(a)] For \(\ell\ge3c+r-1\), setting \(y_\beta=0\) for
\(\beta\ge3c+r-1\) and \(z_\gamma=0\) for \(\gamma\ge3c+2r-2\) restricts
\(C(r,\ell)\) to \(C(r,3c+r-1)\).
\item[(b)] Label the indices of legs \(2\) and \(3\) of \(C(r,3c+r-1)\)
by the intervals
\[
\begin{array}{c|ccc}
 &\mathsf{left}&\mathsf{mid}&\mathsf{right}\\ \hline
 \beta\text{ (leg 2)}&[0,c-1]&[c,2c+r-2]&[2c+r-1,3c+r-2]\\
 \gamma\text{ (leg 3)}&[0,c+r-2]&[c+r-1,2c+r-2]&[2c+r-1,3c+2r-3]
\end{array}
\]
In every term \(x_\alpha y_\beta z_{\alpha+\beta}\), either \(\beta\) and
\(\gamma=\alpha+\beta\) have the same label, or \(\beta\) is
\(\mathsf{mid}\) and \(\gamma\) is \(\mathsf{left}\) or \(\mathsf{right}\).
The terms with both labels \(\mathsf{left}\) form \(C(r,c)\); those with
both labels \(\mathsf{right}\) form \(C(r,c)\) after subtracting
\(2c+r-1\) from \(\beta\) and \(\gamma\); and those with both labels
\(\mathsf{mid}\) form \(C(r,c)^{(23)}\) after \(\alpha\mapsto r-1-\alpha\),
\(\beta\mapsto\beta-c\) and \(\gamma\mapsto\gamma-(c+r-1)\).
\item[(c)] For \(d\ge2\) and \(1\le c<b\), \(b^+\ge3c+r-1\) if and only if
\((d+2)c\le db-r+1\).
\end{enumerate}
\end{lemma}

\begin{proof}
(a) If \(\beta<3c+r-1\), then \(\alpha+\beta\le3c+2r-3\), so the kept
terms are those of \(C(r,3c+r-1)\).
(b) A \(\mathsf{left}\) \(\beta\) gives \(\gamma\le c+r-2\), which is
\(\mathsf{left}\); a \(\mathsf{right}\) \(\beta\) gives
\(\gamma\ge\beta\), which is \(\mathsf{right}\); and a \(\mathsf{mid}\)
\(\gamma\) gives \(\beta=\gamma-\alpha\in[c,2c+r-2]\), which is
\(\mathsf{mid}\). So different labels force a \(\mathsf{mid}\) \(\beta\)
and an outer \(\gamma\). The \(\mathsf{left}\) part consists of the terms
with \(\beta<c\), and the \(\mathsf{right}\) part is a translate. The
\(\mathsf{mid}\) part consists of the terms with \(\gamma\) \(\mathsf{mid}\);
with \(\alpha'=r-1-\alpha\) and \(\gamma'=\gamma-(c+r-1)\in[0,c)\) we get
\(\beta=c+\alpha'+\gamma'\), so after the substitutions the part is
\(\sum_{\alpha'<r,\gamma'<c}x_{\alpha'}y_{\alpha'+\gamma'}z_{\gamma'}\),
which is \(C(r,c)\) with legs \(2\) and \(3\) exchanged.
(c) \(b^+-(3c+r-1)=db-(d+2)c-r+1\).
\end{proof}

Lemma~\ref{lem:stage-types} below removes the cross terms of~(b),
together with the third sum of~\eqref{eq:stage-tensor}.

\subsection{The type restriction of a stage}
\label{subsec:stage-types}

Fix a stage \((r,b)\to(c,r)\) of arity \(d\). Each coordinate on legs
\(2\) and \(3\) of \(U'\) gets a \emph{letter}: \(y^{(j)}_i\) and
\(z^{(j)}_i\) get \(\mathsf K_j\); \(y^{\mathrm q}_\beta\) gets its
leg-\(2\) label from Lemma~\ref{lem:reset}(b) if \(\beta<3c+r-1\), and
\(\mathsf{cropped}\) otherwise; \(z^{\mathrm q}_\gamma\) gets its leg-\(3\)
label if \(\gamma<3c+2r-2\), and \(\mathsf{cropped}\) otherwise. Since
\(b^+\ge3c+r-1\), the quotient coordinates that are not
\(\mathsf{cropped}\) carry \(C(r,3c+r-1)\) and every label occurs. The other letters form three \emph{classes}: \emph{kernel}
(\(\mathsf K_0,\ldots,\mathsf K_{d-2}\)), of \emph{weight}~\(2\);
\emph{outer} (\(\mathsf{left},\mathsf{right}\)), of weight~\(1\); and
\emph{middle} (\(\mathsf{mid}\)), of weight~\(0\). The \emph{diagonal
part} \(U'_\chi\) of such a letter \(\chi\) is the sum of the terms of
\(U'\) with letter \(\chi\) on both legs. By
Lemmas~\ref{lem:shifted-chain}(c) and~\ref{lem:reset}, coordinate
permutations of each leg identify \(U'_{\mathsf K_j}\),
\(U'_{\mathsf{left}}\) and \(U'_{\mathsf{right}}\) with \(C(r,c)\), and
\(U'_{\mathsf{mid}}\) with \(C(r,c)^{(23)}\), the last one reversing the
first index.

\begin{lemma}[Stage type restriction]\label{lem:stage-types}
Let \((r,b)\to(c,r)\) be a stage of arity \(d\), let \(N\ge1\), and let
\(\tau=(\nK,\nO,\nmid)\) be nonnegative integers with sum \(N\), a
\emph{type of width \(N\)}. In \(U'^{\otimes N}\), with product
coordinates on legs \(2\) and \(3\), the \emph{word} of a coordinate is
its sequence of letters, and its \emph{type} is the triple of the
numbers of its kernel, outer and middle letters. After the substitutions
of Lemma~\ref{lem:shifted-chain}(c) in every position, setting to zero
every coordinate on legs \(2\) and \(3\) whose type is not \(\tau\) gives
a restriction of \(U^{\otimes N}\), with integer coefficients, to a sum
\(T=\sum_wT_w\) of
\begin{equation}\label{eq:stage-sector-count}
 M=\frac{N!}{\nK!\,\nO!\,\nmid!}\,(d-1)^{\nK}\,2^{\nO}
\end{equation}
sectors sharing the first leg \(X=(\mathbb F^r)^{\otimes N}\), one for each
word \(w\) of type \(\tau\) (a word that contains \(\mathsf{cropped}\) has a
type of width less than \(N\)). The sector
\(T_w=U'_{w_1}\otimes\cdots\otimes U'_{w_N}\) consists of the terms with
word \(w\) on both legs. Moving the middle positions to the end, keeping
the order within each group, and reversing the first index at middle
positions identifies \(T_w\) with
\begin{equation}\label{eq:stage-sector-tensor}
 \mathcal V=C(r,c)^{\otimes(N-\nmid)}\bigl(C(r,c)^{(23)}\bigr)^{\otimes\nmid}
\end{equation}
by coordinate permutations of each leg; on the first leg they depend
on \(w\).
\end{lemma}

\begin{proof}
\emph{One position.} Let \(x_\alpha y'_{\mathbf u}z'_{\mathbf v}\) be a
term of \(U'\) with letters \(\chi\ne\mathsf{cropped}\) on leg \(2\) and
\(\chi'\ne\mathsf{cropped}\) on leg \(3\). We claim that \(\chi'=\chi\) or
that \(\chi'\) has larger weight. By~\eqref{eq:stage-tensor}, a kernel
coordinate \(y^{(j)}_i\) occurs only with \(z^{(j)}_{i+\alpha}\). A
quotient coordinate, of weight at most \(1\), occurs either with a kernel
coordinate, of weight \(2\), or in a term
\(x_\alpha y^{\mathrm q}_\beta z^{\mathrm q}_{\beta+\alpha}\) of
\(C(r,b^+)\) with \(\beta<3c+r-1\), which by Lemma~\ref{lem:reset} has
equal labels or goes from \(\mathsf{mid}\) to an outer label.

\emph{The power.} Setting coordinates to zero is a restriction, so
\(U^{\otimes N}\restr U'^{\otimes N}\restr T\). The coefficient of a term
of \(U'^{\otimes N}\) is the product of the coefficients in the \(N\)
positions. In a surviving term with words \(w\) on leg \(2\) and \(w'\)
on leg \(3\), the weight of \(w'_q\) is therefore at least that of
\(w_q\) at every position \(q\). Both words have type \(\tau\)
and hence weight \(2\nK+\nO\), so the weights agree at every position,
and by the claim \(w'=w\). Thus \(T\) is the sum of the \(T_w\), which use
disjoint coordinates on legs \(2\) and \(3\), and
\(T_w=\bigotimes_qU'_{w_q}\).

\emph{Count.} A word of type \(\tau\) is given by the positions of the
three classes, a kernel letter at each kernel position and an outer
letter at each outer position, which gives~\eqref{eq:stage-sector-count}.
The identification follows from those of the parts \(U'_\chi\).
\end{proof}

In Example~\ref{ex:shifted-chain}, where~\eqref{eq:stage-conditions}
reads \(4\le5\), the crop removes \(y^{\mathrm q}_4\) and
\(z^{\mathrm q}_5\), and the cross terms are
\(x_0y^{\mathrm q}_1z^{\mathrm q}_1\) (\(\mathsf{mid}\) to
\(\mathsf{left}\)), \(x_1y^{\mathrm q}_2z^{\mathrm q}_3\) (\(\mathsf{mid}\)
to \(\mathsf{right}\)) and \(x_0y^{\mathrm q}_3z^{(0)}_1\)
(\(\mathsf{right}\) to \(\mathsf K_0\)). Fixing only the number of middle
letters, as in \citep[Lemma~4.5]{FinCon26}, would keep the last one.

\subsection{The stage relation}
\label{subsec:stage-relation}

\begin{proposition}[One stage]\label{prop:stage}
Let \((r,b)\to(c,r)\) be a stage of arity \(d\), let \(N\ge1\), let
\(\tau\) be a type of width \(N\), and let \(M\) be given
by~\eqref{eq:stage-sector-count}. There is a degeneration over \(\Q\)
\begin{equation}\label{eq:stage-relation}
 (2M-1)^6\,Q_{r,b}^{\otimes N}\,\MM_{d^{2N}}\degen M^6\,Q_{c,r}^{\otimes N}\,\MM_{M^2}
\end{equation}
of degree at most \(6(M-1)\). In each orientation \(\pi\in S_3\), its maps
are the restriction of Lemma~\ref{lem:stage-types}, the
separator~\eqref{eq:prelim-separator} with shared leg \(\pi(1)\), and
coordinate permutations that do not depend on \(t\).
\end{proposition}

\begin{proof}
Let \(U\) be the tensor~\eqref{eq:stage-U}; since
\(B_X(d)^{\otimes N}\cong B_X(d^N)\),
\(U^{\otimes N}\cong C(r,b)^{\otimes N}B_X(d^N)\). Apply the restriction
of Lemma~\ref{lem:stage-types} to each of \(2M-1\) copies of
\(U^{\otimes N}\), then the separator~\eqref{eq:prelim-separator}, of
degree at most \(M-1\), and then the reindexing of
Section~\ref{subsec:prelim-sectors} built from the identifications
\(T_w\cong\mathcal V\). The first and last maps do not depend on \(t\),
and the last is defined on all target variables. By
\citep[Lemma~2.1]{FinCon26}, for every \(\pi\in S_3\) this gives
\((2M-1)(U^{\otimes N})^\pi\degen M(\mathcal VB_X(M))^\pi\), of degree at
most \(M-1\), with shared leg \(\pi(1)\), and tensoring the six
orientations with one parameter \(t\) gives
\[
 (2M-1)^6\bigotimes_{\pi\in S_3}(U^{\otimes N})^\pi
 \degen M^6\bigotimes_{\pi\in S_3}\bigl(\mathcal VB_X(M)\bigr)^\pi
\]
of degree at most \(6(M-1)\). All maps are rational. It remains to
identify both sides by coordinate permutations within each global leg,
which change neither the degree nor the constant term.
By~\eqref{eq:prelim-mm}, \(\bigotimes_\pi(U^{\otimes N})^\pi\cong
Q_{r,b}^{\otimes N}\bigotimes_\pi B_X(d^N)^\pi\cong Q_{r,b}^{\otimes N}\MM_{d^{2N}}\)
and \(\bigotimes_\pi B_X(M)^\pi\cong\MM_{M^2}\). Since
\((T^{(23)})^\pi=T^{\pi(23)}\) for every tensor \(T\),
\[
 \bigotimes_{\pi\in S_3}\mathcal V^\pi
 =\bigotimes_{\pi\in S_3}\bigl(C(r,c)^\pi\bigr)^{\otimes(N-\nmid)}
 \bigl(C(r,c)^{\pi(23)}\bigr)^{\otimes\nmid},
\]
and as \(\pi\) runs through \(S_3\), so does \(\pi(23)\). Hence
\(\bigotimes_\pi\mathcal V^\pi\cong Q_{r,c}^{\otimes N}\cong Q_{c,r}^{\otimes N}\)
by~\eqref{eq:prelim-Q}.
\end{proof}

\begin{corollary}[Exact stage]\label{cor:exact-stage}
In the setting of Proposition~\ref{prop:stage}, let \(k\ge1\) and
\(E=1+6k(M-1)\). Then, over \(\Q\),
\begin{equation}\label{eq:stage-exact}
 E\,(2M-1)^{6k}\,Q_{r,b}^{\otimes Nk}\,\MM_{d^{2Nk}}
 \restr M^{6k}\,Q_{c,r}^{\otimes Nk}\,\MM_{M^{2k}} .
\end{equation}
\end{corollary}

\begin{proof}
Apply exactification in its rational form \citep[Lemmas~2.2
and~9.1(c)]{FinCon26} to~\eqref{eq:stage-relation}, with degree bound
\(\theta=6(M-1)\) and \(k\theta+1=E\), and use \((mT)^{\otimes k}\cong m^kT^{\otimes k}\) and
\(\MM_a^{\otimes k}\cong\MM_{a^k}\).
\end{proof}

\begin{remark}[Both signs]\label{rem:both-signs}
For one class \(\mathcal A_\sigma\), the proof gives
\(\bigotimes_{\pi\in\mathcal A_\sigma}\mathcal V^\pi\cong
(S^\sigma_{r,c})^{\otimes(N-\nmid)}(S^{-\sigma}_{r,c})^{\otimes\nmid}\),
since \(\pi(23)\in\mathcal A_{-\sigma}\). For \(c=r\),
\citep[(4.6)]{FinCon26} closes this with \(S^+_{r,r}\cong S^-_{r,r}\). For
\(r\ne c\) with \(r,c\ge2\), the identities of \citep[(2.6)]{FinCon26} do
not relate \(S^+_{r,c}\) to \(S^-_{r,c}\), so we use all six
orientations, as in the proof of \citep[Lemma~4.2]{OAI26}. This costs
nothing in the exponent. A threefold relation of degree \(3(M-1)\), if it
existed, would have the auxiliaries
\(\bigotimes_{\pi\in\mathcal A_\sigma}B_X(e)^\pi\cong\MM_e\)
\citep[(2.2)]{FinCon26} with \(e=d^N\) and \(e=M\). At power \(2k\) it
would give the same copy counts, auxiliary sides and number
\(E=1+2k\cdot3(M-1)\) of exactification copies
as~\eqref{eq:stage-exact}; only its odd powers, which give smaller sides
but no smaller exponent, have no sixfold counterpart (compare
\citep[Appendix~A.4]{FinCon26}).
\end{remark}

\subsection{Chains}
\label{subsec:chain-counts}

A \emph{chain} of \(J\ge1\) stages from the \emph{source pair}
\((r_1,b_1)\) to the \emph{terminal pair} \((1,t)\), \(t\ge1\), consists
of stages \((r_i,b_i)\to(c_i,r_i)\) of arities \(d_i\), \(1\le i\le J\),
with \((r_{i+1},b_{i+1})=(c_i,r_i)\) for \(i<J\) and \((c_J,r_J)=(1,t)\).
Its \emph{evaluation size} is \(s=r_1+b_1-1\).

\begin{proposition}[Chain counts]\label{prop:chain-counts}
Fix a chain as above and integers \(N,k\ge1\). At each stage \(i\) fix a
type of width \(N\), let \(M_i\) be its sector
count~\eqref{eq:stage-sector-count} with \(d=d_i\), and put
\(E_i=1+6k(M_i-1)\). Then \(r_*\MM_A\restr m_*\MM_B\) over \(\Q\), where
\begin{equation}\label{eq:stage-chain-counts}
\begin{aligned}
 A&=\prod_{i=1}^Jd_i^{2Nk},&\qquad
 B&=t^{2Nk}\prod_{i=1}^JM_i^{2k},\\
 r_*&=s^{6Nk}\prod_{i=1}^JE_i\,(2M_i-1)^{6k},&\qquad
 m_*&=\prod_{i=1}^JM_i^{6k}.
\end{aligned}
\end{equation}
If \(B>A\), this is a relative scheme \citep[Definition~2.3]{FinCon26}, its
pooled exponent \(p=\log(r_*/m_*)/\log(B/A)\) satisfies \(p\ge2\), and
\(\omega\le p\) over every field of characteristic zero
\citep[Lemma~2.4]{FinCon26}.
\end{proposition}

\begin{proof}
Put \(\mathcal T_i=Q_{r_i,b_i}^{\otimes Nk}\) for \(1\le i\le J+1\), with
\((r_{J+1},b_{J+1})=(1,t)\), so that \(\mathcal T_{J+1}\cong\MM_{t^{2Nk}}\)
by~\eqref{eq:prelim-Q} and~\eqref{eq:prelim-mm}. Put
\(\rho_i=E_i(2M_i-1)^{6k}\) and \(\rho'_i=M_i^{6k}\). Since
\(Q_{c_i,r_i}=Q_{r_{i+1},b_{i+1}}\), Corollary~\ref{cor:exact-stage} reads
\(\rho_i\mathcal T_i\MM_{d_i^{2Nk}}\restr\rho'_i\mathcal T_{i+1}\MM_{M_i^{2k}}\). With
\(A_{\ge j}=\prod_{i\ge j}d_i^{2Nk}\) and \(B_{<j}=\prod_{i<j}M_i^{2k}\),
we show by induction on \(j=1,\ldots,J+1\) that
\[
 \Bigl(\prod_{i=1}^J\rho_i\Bigr)\mathcal T_1\MM_A\restr
 \Bigl(\prod_{i<j}\rho'_i\Bigr)\Bigl(\prod_{i\ge j}\rho_i\Bigr)
 \mathcal T_j\,\MM_{A_{\ge j}B_{<j}} .
\]
For \(j=1\) the two sides are equal. For the step, write
\(\MM_{A_{\ge j}B_{<j}}\cong\MM_{d_j^{2Nk}}\MM_{A_{\ge j+1}B_{<j}}\). The
right side is a direct sum of copies of
\(\rho_j\mathcal T_j\MM_{d_j^{2Nk}}\MM_{A_{\ge j+1}B_{<j}}\). Applying
Corollary~\ref{cor:exact-stage} for stage \(j\), tensored with
\(\MM_{A_{\ge j+1}B_{<j}}\), to each copy turns it into
\(\rho'_j\mathcal T_{j+1}\MM_{A_{\ge j+1}B_{<j+1}}\), which is the claim for
\(j+1\). For \(j=J+1\), the right side is isomorphic to \(m_*\MM_B\).
Finally, \eqref{eq:prelim-Q} gives \(s^{6Nk}\MM_1\restr \mathcal T_1\) over \(\Q\),
so \(r_*\MM_A\restr(\prod_i\rho_i)\mathcal T_1\MM_A\), and composing gives
\(r_*\MM_A\restr m_*\MM_B\). The rest is
Section~\ref{subsec:prelim-schemes}.
\end{proof}

\begin{remark}[Stage-dependent \(N\) and \(k\)]\label{rem:varying-powers}
The proof needs only that the \(Q\)-power in the target of each stage is
the \(Q\)-power in the source of the next. So stage \(i\) may use its own
\(N_i,k_i\ge1\), a type of width \(N_i\) and \(E_i=1+6k_i(M_i-1)\),
provided \(N_ik_i=N_1k_1\) for all \(i\).
The same induction then gives \(A=\prod_id_i^{2N_1k_1}\),
\(B=t^{2N_1k_1}\prod_iM_i^{2k_i}\),
\(r_*=s^{6N_1k_1}\prod_iE_i(2M_i-1)^{6k_i}\) and
\(m_*=\prod_iM_i^{6k_i}\).
\end{remark}

\begin{remark}[Fields]\label{rem:field-scope}
Lemmas~\ref{lem:shifted-chain}, \ref{lem:reset}
and~\ref{lem:stage-types} use integer matrices with integer inverses and
coordinate selections, so they hold over every field.
Proposition~\ref{prop:stage}, Corollary~\ref{cor:exact-stage} and
Proposition~\ref{prop:chain-counts} are proved over \(\Q\), hence hold over
every field of characteristic zero; their coefficients come from the
interpolation of \citep[Lemma~9.1]{FinCon26} at the nodes
\(1,\ldots,s\) and, at each stage \(i\), \(1,\ldots,M_i\),
\(1,\ldots,2M_i-1\) and \(1,\ldots,E_i\). We do not treat positive
characteristic (Section~\ref{sec:open}).
\end{remark}
